\documentclass[11pt,a4paper]{article}
\usepackage[margin=24mm]{geometry}
\usepackage{amsmath,amssymb,hyperref}
\hypersetup{hidelinks}
\title{Stanley Supplementary Exercise 88\\The power-sum image of the hook-content transformation}
\author{Statement: Richard P. Stanley\\Paired original proof: Richard P. Stanley}
\date{Statement version 4 June 2023; author proof arXiv:0807.0383v3}
\begin{document}
\maketitle
\section*{Original exercise}
Define a \(\mathbb{Q}\)-linear transformation \(\varphi:\Lambda_{\mathbb{Q}}\to\mathbb{Q}[t]\) by
\[
 \varphi(s_\lambda)=\frac{\prod_{i=1}^{n}(t+\lambda_i+n-i)}{H_\lambda},
\]
where \(\lambda=(\lambda_1,\dots,\lambda_n)\vdash n\) and \(H_\lambda\) denotes the product of the hook lengths of \(\lambda\).
Show that for any \(\mu\vdash n\) with \(\ell(\mu)=\ell\) and \(m_1(\mu)=m\) (the number of parts of \(\mu\) equal to \(1\)), we have
\[
 \varphi(p_\mu)=(-1)^{n-\ell}\sum_{i=0}^{m}\binom{m}{i}t(t+1)\cdots(t+i-1).
\]

\section*{Paired author proof}
The official solutions explicitly cite Lemma 3.2 of R.~Stanley, \emph{Some combinatorial
properties of hook lengths, contents, and parts of partitions}, Ramanujan J.~23 (2010),
91--105; arXiv:0807.0383v3. The lemma as printed there is exactly the displayed identity,
with the paper's variable \(v\) in place of \(t\) (the exercise writes \(\varphi:\Lambda_{\mathbb{Q}}\to\mathbb{Q}[t]\), the paper writes \(\varphi:R\to\mathbb{Q}[v_1,\dots,v_k]\) and sets \(v=v_1\)).
The complete native \LaTeX{} source and the author PDF of version 3 are retained with this statement.
The paired source uses the paper's definition of \(\varphi\) (Section~3) and the full proof of
Lemma~3.2: expanding \(e_{n-i}\) in power sums gives the coefficient of \(p_\mu\) in equation~(16),
and applying \(\varphi\) to the Cauchy identity~(17) followed by linear independence of the
\(p_\mu\) closes the argument. Lemma~3.1 supplies the expansion used at equation~(11).

\section*{Sources}
Original statement: Stanley, exercise 88, physical page 27,
\url{https://math.mit.edu/~rstan/ec/ch7supp.pdf}.
Official solution pointer: version 24 March 2023, physical page 40,
\url{https://math.mit.edu/~rstan/ec/ch7suppsol.pdf}.
Original author proof: \url{https://arxiv.org/abs/0807.0383v3}.
\end{document}
